% Created 2026-09-29 mar. 17:57
% Intended LaTeX compiler: pdflatex
\documentclass[11pt]{article}
\usepackage[utf8]{inputenc}
\usepackage[T1]{fontenc}
\usepackage{graphicx}
\usepackage{longtable}
\usepackage{wrapfig}
\usepackage{rotating}
\usepackage[normalem]{ulem}
\usepackage{amsmath}
\usepackage{amssymb}
\usepackage{capt-of}
\usepackage{hyperref}
\author{Romain Capliez}
\date{\today}
\title{}
\hypersetup{
 pdfauthor={Romain Capliez},
 pdftitle={},
 pdfkeywords={},
 pdfsubject={},
 pdfcreator={Emacs 29.1 (Org mode 9.6.6)}, 
 pdflang={English}}
\begin{document}

\tableofcontents

\section{Objectif}
\label{sec:org1be0248}

Etudier les impacts macro d'un choc sur les prix des métaux critiques et comparer ces effets à ceux d'un choc sur les prix du pétrole.

Les effets sont-ils les mêmes ?
Certains pays sont-ils plus impactés que d'autres ?

\section{Variables endogènes :}
\label{sec:orgc89b3a3}
\begin{itemize}
\item Energy consummer prices
\item Core consummer prices
\item Industrial production
\item Stock prices
\item Short term interest rate
\item Long-term interest rate
\item Price of electricity ?
\item Industrial production price (pass through to consummer price?.)
\end{itemize}

\section{Variables exogènes:}
\label{sec:orgddf725b}
\begin{itemize}
\item Prix du cuivre (autres métaux du LME ?) (comment faire exogénéité ? -> Alessandri2025 ?)
\item Prix du pétrole (ou chocs pourr exogénéité)
\end{itemize}

\section{Variables de transition:}
\label{sec:orgbc7ab05}
\begin{itemize}
\item Prix du pétrole : effet différent selon crise ou pas
\item Intensité métalique : Effet différents pour les pays + dépendants aux métaux
\item Temps : Effet qui évolue dans le temps
\end{itemize}

\section{Méthodologie}
\label{sec:orgcbd8db5}
Projections locales panel non-linéaire à transition lisse : 
\begin{itemize}
\item Permet d'avoir un impact global mais également de regarder pays par pays
\item Intégration "facile" de la non-linéarité
\item Le fait que les prix du pétrole et des métaux soient communs à tous les pays est-il un problème? Se confond avec effets fixes temporels
\end{itemize}


Pass-through methodology : Alessandri2025 ?

\begin{equation}
\label{eq:1}
PT_h^s = \frac{\sum_{1:h} IRF_h^s \left( \text{core hicp} \right)}{\sum_{1:h} IRF_h^s \left( \text{energy hicp} \right) }
\end{equation}


What is the core inflation cost of a 1\% rise in energy prices? And does that depend on whether the rise in energy prices originates in the gas or the oil market?
\end{document}